\subsection{Checks for \texorpdfstring{Section~\ref{sec:evo}}{evolutionary individuality}}\label{chk:evo}

The checks below evaluate the closed forms of \cref{sec:evo} on random parameter draws, compare the transmission matrix with a simulation of the fission life cycle at the level of individual members, and iterate the selection recursions against their fixed points and against a simulated population. Every statement of \cref{sec:evo,app:evo} has a proof; the checks guard the algebra and the modeling of the life cycle, and each one ends by naming what it does not test. Deviations are the largest found over the stated draws.

% script: evo_heritability.py, function check_identity (seed 1)
\begin{numcheck}[Identity rate of the manuscript re-derived]\label{chk:evo-identity}
The section builds on \cref{prop:agy-identity}, so its formulas were re-derived and tested independently of \cref{chk:agy-identity}. On 20000 draws with $\varepsilon$, $\beta$, $g$, $\omega$ and $\delta$ uniform on $[0,3]$, $\omega$ floored at $10^{-3}$, and $\mathsf f$ uniform on $[0,1]$, all of which had $\det\mathsf A>10^{-9}$, the closed form of $\varpi_\pm$ agreed with the numerically computed eigenvalues of $-\mathsf A$ to $3.6\times10^{-15}$. The rate $\varpi_-$ never exceeded the comparison value $\delta+\omega\tilde\varepsilon/\mathsf p_1$, the smallest gap being $6.3\times10^{-7}$, and on the 6163 draws with $\mathsf p_2\le\frac12\mathsf p_1$ it fell below that value by at most $0.794$ times $2\omega\mathsf p_2/\mathsf p_1$. With $g=1$ and $\varepsilon=0.01$ the identity time without a record was $0.99$, $1.96$, $9.09$ and $50.0$ residence times at $\mathsf f=0$, $0.5$, $0.9$ and $0.99$, and with a record ($\beta=100$, $\omega=0.01$, $\mathsf f=0.9$, $\delta=0.001$) it was $989.3$. These agree with the values reported in \cref{chk:agy-identity}, and no discrepancy with \cref{prop:agy-identity}(a) was found. Part (b) of that proposition was not tested numerically.
\end{numcheck}

% script: evo_heritability.py, functions check_heritability (seed 2) and check_read_rate (seed 5)
\begin{numcheck}[Heritability closed forms and the record channel]\label{chk:evo-heritability}
On 40000 draws with $\varepsilon$, $\beta$, $g$, $\omega$ and $\delta$ uniform on $[0,2]$, $\omega$ floored at $10^{-3}$, $\mathsf f$, $\mathsf c$ and $\bar{\mathsf b}$ uniform on $[0,1]$ and $\mathsf T$ uniform on $[0.05,6]$, none of which had $\tilde\beta<10^{-3}$, the closed form \eqref{eq:evo-semigroup} agreed with $e^{\mathsf T\mathsf A}$ computed by eigendecomposition to $1.9\times10^{-15}$, and \eqref{eq:evo-heritability} agreed with the numerically computed spectral radius of $\mathsf M$ to $2.2\times10^{-16}$. Every draw satisfied the three inequalities of \eqref{eq:evo-sandwich}: the smallest values of $\mathsf h-e^{-\varpi_-\mathsf T}\min\{\bar{\mathsf b},\mathsf c\mathsf k\}$, of $e^{-\varpi_-\mathsf T}\max\{\bar{\mathsf b},\mathsf c\mathsf k\}-\mathsf h$ and of $\mathsf h-e^{-\varpi_-\mathsf T}\max\{\mathsf c\mathsf u,\bar{\mathsf b}(1-\mathsf u)\}$ were $5.6\times10^{-10}$, $7.9\times10^{-10}$ and $4.8\times10^{-11}$, all nonnegative, and $1-\mathsf h\ge0.048$ throughout. The equivalence of $\varpi_-<\tilde\varepsilon$ and $\delta<\tilde\varepsilon$ in \cref{thm:evo-record-channel}(b) held on every draw. On the draws with $\delta\ge\tilde\varepsilon$ the ratio of $\mathsf h/\mathsf h_0$ to $\max\{\bar{\mathsf b},\mathsf c\mathsf k\}/(\bar{\mathsf b}+(1-\bar{\mathsf b})\mathsf f)$ was at most $0.9991$, a consequence of part (d); the bound of part (d) itself is the upper inequality of \eqref{eq:evo-sandwich} divided by $\mathsf h_0$. For \cref{ex:evo-worked} the script returned $\tau_{\mathrm{id}}=60.27$, $1/\tilde\varepsilon=0.990$, $\mathsf u=0.9890$, the limit $0.9406$ and the entries of \cref{tab:evo-worked}, with $\mathsf he^{\mathsf T/\tau_{\mathrm{id}}}=0.9410$, $0.9406$, $0.9406$ and $0.9406$ at $\mathsf T=1$, $5$, $20$ and $60$.

For the dependence on the read rate discussed in \cref{rem:evo-bourrat-detail}, on 20000 further draws from the same law with $\beta$ floored at $10^{-2}$, none of which had $\varpi_+-\varpi_-<0.05$, the formula $\partial\varpi_-/\partial\beta=(\delta-\varpi_-)/(\varpi_+-\varpi_-)$ agreed with a central difference of step $10^{-5}$ to $4.5\times10^{-8}$, and its sign was negative exactly on the draws with $\delta<\tilde\varepsilon$. On 2094 of the 13853 draws with $\delta<\tilde\varepsilon$ the heritability at the drawn generation time was nevertheless lower at $\beta+10^{-3}$ than at $\beta$. One instance with a record that is blank at birth, $\varepsilon=0.1$, $g=1$, $\mathsf f=0.5$, $\omega=0.2$, $\delta=0$, $\mathsf c=0$, $\bar{\mathsf b}=0.5$ and $\mathsf T=0.2$: at $\beta=0.5$, $1$ and $2$ the identity time is $14.42$, $18.72$ and $27.19$ while $\mathsf h=0.3647$, $0.3308$ and $0.2723$. Heritability at a fixed generation time is therefore not increasing in the read rate, although the identity time is whenever $\delta<\tilde\varepsilon$.

Not tested: \cref{thm:evo-record-channel}(e) was not simulated by this script, and only its coefficient $\bar{\mathsf b}e^{-\mathsf p_1\mathsf T}$ is tabulated (its conditional mean is computed exactly in \cref{chk:evo-independent}); the limits $\mathsf T\to\infty$ and $\mathsf T\to0$ of part (a) are covered by \cref{chk:evo-bottleneck}.
\end{numcheck}

% script: evo_heritability.py, function check_bottleneck (seed 3)
\begin{numcheck}[Founders against recruits]\label{chk:evo-bottleneck}
On 20000 draws with rates uniform on $[0,2]$, $\omega$ floored at $10^{-3}$, and $\mathsf f$ and $\mathsf c$ uniform on $[0,1]$, the bracket $\min\{\delta,\tilde\varepsilon\}\le\varpi_-\le\max\{\delta,\tilde\varepsilon\}$ of \cref{lem:evo-semigroup}(v) held on every draw with tolerance $10^{-12}$. The long-generation prefactor was computed as the spectral radius of $e^{\varpi_-\mathsf T}\mathsf G(\mathsf T)\mathsf F$ at $(\varpi_+-\varpi_-)\mathsf T=60$ for $\bar{\mathsf b}\in\{0,\frac12,1\}$; it agreed with $\mathsf P(\bar{\mathsf b})=(1-\mathsf u)\bar{\mathsf b}+\mathsf u\mathsf c\mathsf k$ to $6.7\times10^{-16}$, the difference $\mathsf P(1)-\mathsf P(0)$ agreed with the slope \eqref{eq:evo-bottleneck-slope} to $1.7\times10^{-15}$, and the sign of the difference agreed with the sign of $\omega(1-\mathsf c\mathsf f)+\delta-\varpi_-$ on every draw. The spectral radius of $\mathsf F$ equaled $\max\{\bar{\mathsf b},\mathsf c\}$ to $1.1\times10^{-16}$ at one uniformly drawn $\bar{\mathsf b}$ per draw, which is the limit $\mathsf T\to0$ of \cref{thm:evo-record-channel}(a). An instance of the inversion, with faithful copy and faithful recruits and a record that does not decay, $\varepsilon=0.2$, $\beta=1$, $g=1$, $\mathsf f=1$, $\omega=0.5$, $\delta=0$, $\mathsf c=1$ and $n=10$, has $\varpi_-=0.0376>\delta+\omega(1-\mathsf c\mathsf f)=0$, and the heritability at $\mathsf b=0$, $1$, $5$ and $10$ is $0.99166$, $0.99089$, $0.98732$ and $0.98140$ at $\mathsf T=0.5$, $0.84064$, $0.83946$, $0.83471$ and $0.82878$ at $\mathsf T=5$, and $0.47856$, $0.47788$, $0.47518$ and $0.47180$ at $\mathsf T=20$. In this instance every founder lowers the heritability at all three generation times, including the short one, which \cref{prop:evo-bottleneck} does not cover; the threshold $\mathsf T_0$ of that proposition was not computed, and monotonicity in $\mathsf b$ at intermediate generation times was not tested on the random draws.
\end{numcheck}

% script: evo_fission_mc.py (seed 7)
\begin{numcheck}[Fission life cycle at the level of members]\label{chk:evo-fission}
One generation of \cref{def:evo-fission} was simulated member by member for 20000 parents with one daughter each, with $n=8$, $\mathsf b=2$, $\varepsilon=0.05$, $\beta=0.8$, $g=0.5$, $\mathsf f=0.6$, $\omega=0.4$, $\delta=0.02$, $\mathsf c=0.9$, $\mathsf T=2$, $\mathsf s_\nu=0.3$ and $\mathsf s_z=0.2$. Parents' member states were a common standard normal component plus independent deviations of standard deviation $0.5$, and the parent's record was $0.7$ times the common component plus noise of standard deviation $0.6$. The relaxation used an Euler step of $0.01$ (200 steps) with mean-conserving consensus among members at rate $1$, private diffusion of intensity $0.1$, and replacement of each member in each step with probability $1-e^{-0.005}$. The conditional mean of this discrete scheme is exactly linear with a known $2\times2$ one-step matrix, so the Monte Carlo error and the discretization error can be separated. The regression of the daughter's state at birth on the parent's state returned $(0.2473,0.4061;-0.0019,0.8988)$, rows separated by a semicolon, against $\mathsf F=(0.25,0.405;0,0.9)$. The regression of the daughter's final state returned $(0.0502,0.5910;0.0496,0.6980)$ against $(0.0525,0.5901;0.0513,0.6984)$ for the discrete scheme, a largest deviation of $2.3\times10^{-3}$, and the discrete matrix differs from $\mathsf M=e^{\mathsf T\mathsf A}\mathsf F=(0.0525,0.5904;0.0511,0.6989)$ by at most $4.6\times10^{-4}$. The cross-covariance of final and parental states differed from the product of the discrete matrix with the parental covariance by at most $2.0\times10^{-3}$, for entries up to $0.625$. The parent-offspring regression slope of the heritable trait was $0.7405$ with standard error $0.0012$, against the spectral radius $0.7423$ of the discrete matrix ($1.5$ standard errors) and $\mathsf h=0.7426$. The empirical covariance of the birth state about $\mathsf F\mathsf x$ was $(0.02267,0.01822;0.01822,0.04001)$ against $(0.02239,0.01800;0.01800,0.04000)$ from \eqref{eq:evo-noise} with $\mathsf v_{\mathrm w}$ replaced by its mean over parents; the nominal relative sampling error of a variance estimated from 20000 normal draws is 1 percent. In the variant without a record the slope was $0.4272$ with standard error $0.0013$, against $0.4248$ for the discrete scheme ($1.8$ standard errors) and $\mathsf h_0=0.4246$, and the variance of the members' mean at birth about its conditional mean was $0.05342$ against $0.05422$ from \cref{prop:evo-variance}(c). Not tested: the covariance $\mathsf N_{\mathrm{rel}}$ added during relaxation, which is an input and was not compared with any formula, the convergence along lines of descent in \cref{prop:evo-transmission}(c), and the variance recursion of \cref{prop:evo-variance}(b) over several generations.
\end{numcheck}

% scripts: evo_balance.py, function price_check (seed 11); evo_fission_mc.py, Price block (seed 7)
\begin{numcheck}[Price decomposition]\label{chk:evo-price}
For 500 parents with traits drawn from $\mathcal N(0.5,1)$, numbers of daughters Poisson with mean $e^{0.5(y_j-\bar y)}$, and daughters' traits $\mathsf hy_j$ plus independent normal noise with $\mathsf h=0.8$ and standard deviation $0.3$, the two sides of \eqref{eq:evo-price-identity} differed by at most $4.4\times10^{-16}$ over 400 independent draws of the offspring generation. With the parents and their fitnesses held fixed, the mean realized change of the trait mean over those 400 draws was $0.2731$ with standard error $0.0006$, against $\mathsf h\mathsf d_{\mathrm{sel}}-(1-\mathsf h)\bar y=0.2741$ from \eqref{eq:evo-price} ($1.7$ standard errors); over 20000 further draws it was $0.27419$ with standard error $0.00009$, against $0.27407$ ($1.3$ standard errors). In the member-level simulation of \cref{chk:evo-fission}, with numbers of daughters Poisson with mean $e^{0.4(y_j-\bar y)}$ and the one simulated daughter of each parent weighted by $\mathsf W_j$, the realized change of the mean heritable trait was $0.2780$ against $0.2792$ predicted with the spectral radius of the discrete matrix, with a Monte Carlo standard error of $0.0017$, selection differential $\mathsf d_{\mathrm{sel}}=0.3749$ and $\bar y=-0.0036$. The weighting reproduces the conditional expectation in \eqref{eq:evo-price} and not the sampling variance of a population in which each parent founds $\mathsf W_j$ distinct daughters. The general form \eqref{eq:evo-price-general} for a trait other than the heritable one was not tested separately.
\end{numcheck}

% script: evo_balance.py, functions balance_check (seed 11) and two_type_check
\begin{numcheck}[Selection-transmission balance and the absence of a threshold]\label{chk:evo-balance}
With $\mathsf h=0.85$, $\mathsf s^2=0.04$, $\mathsf v_{\mathrm s}=1$ and $\mathsf y_\ast=2$, 400 iterations of \eqref{eq:evo-balance-recursion} from $(\mathsf V,\bar y)=(0,-3)$ and from $(5,4)$ both gave $(0.113847,0.733532)$, equal to \eqref{eq:evo-balance-fixed} to the six digits printed. A simulated population of 200000 collectives, with parents sampled in proportion to $\mathsf W$ and daughters' traits drawn as in \cref{ass:evo-gaussian}, run for 120 generations from a standard normal start, had variance $0.11386$ (standard deviation across generations $0.00056$) and mean $0.73255$ ($0.00167$) over generations 61 to 120, and mean fitness $0.46069$ ($0.00090$) over generations 61 to 119 against $0.46121$ from \eqref{eq:evo-balance-fitness}. At $\mathsf s^2=0.04$ and $\mathsf v_{\mathrm s}=1$ the retained fraction $\hat y/\mathsf y_\ast$ was $0.0020$, $0.0099$, $0.0292$, $0.0789$, $0.2597$, $0.5228$, $0.7373$, $0.9451$ and $1.0000$ at $\mathsf h=0.05$, $0.2$, $0.4$, $0.6$, $0.8$, $0.9$, $0.95$, $0.99$ and $1$; on a grid of 20001 values of $\mathsf h$ in $[10^{-4},1]$ it was strictly increasing, with smallest value $3.9\times10^{-6}$. For \cref{prop:evo-directional} with $\mathsf r=0.3$ and the same $\mathsf h$ and $\mathsf s^2$, 2000 iterations gave $(0.144144,0.245045)$, equal to \eqref{eq:evo-directional}. For the two-type map \eqref{eq:evo-two-type} with $\mathsf m=3$ and $p_0=\frac12$, the frequency after 20000 generations was $0$ to six decimals at $\mathsf q=0.2$, $2.5\times10^{-5}$ at $\mathsf q=\frac13$, $0.250000$ at $\mathsf q=0.5$ and $0.850000$ at $\mathsf q=0.9$, against the limits $0$, $0$, $0.25$ and $0.85$ of \cref{prop:evo-two-type}; at the threshold $\mathsf m\mathsf q=1$ the decay is algebraic, $p_k=1/(2+2k)$, which is the value found. The simulated population draws daughters from the Gaussian kernel of \cref{ass:evo-gaussian} and therefore tests the recursion under the closure, not the closure itself; the monotonicity of the retained fraction in $\mathsf s^2$ is proved and was not checked numerically by this script; \cref{chk:evo-independent} checks it on a grid.
\end{numcheck}

% script: evo_indep_exact.py (seed 90210), written by the independent checker; no sampling of the fission or relaxation noise
\begin{numcheck}[Independent recomputation by enumeration]\label{chk:evo-independent}
This check was written separately from the scripts behind the six checks above. It uses its own matrix exponential, a Taylor series with scaling and squaring, and it replaces Monte Carlo sampling of the life cycle by exact computation of conditional means and by enumeration of founder sets. On 6000 draws with $\varepsilon$, $\beta$ and $\delta$ equal to $0$ with probability $0.15$ and otherwise uniform on $[0,2]$, $g$ uniform on $[0.05,2]$, $\omega$ uniform on $[10^{-3},2]$, $\mathsf f$ and $\mathsf c$ equal to $0$ and to $1$ with probability $0.15$ each and otherwise uniform on $[0,1]$, $n$ uniform on $\{2,\dots,12\}$, $\mathsf b$ uniform on $\{0,\dots,n\}$ and $\mathsf T$ log-uniform on $[0.01,30]$, retained when $\tilde\beta>10^{-6}$ and $\det\mathsf A>10^{-9}$, the closed form \eqref{eq:evo-semigroup} agreed with the exponential to $1.8\times10^{-13}$ and \eqref{eq:evo-heritability} with the numerically computed spectral radius to $2.2\times10^{-16}$. No draw violated the ordering of the roots, $0<\mathsf u<1$, the positivity and row-sum bound of $\mathsf G(\mathsf T)$, the bracket of \cref{lem:evo-semigroup}(v), $\mathsf h<1$, any of the four inequalities of \eqref{eq:evo-sandwich}, or the ratio bounds of \cref{thm:evo-record-channel}(c) and (d), at relative tolerance $10^{-12}$; the largest heritability was $0.999946$. The limit of part (a) as $\mathsf T\to\infty$ was reproduced to $4.2\times10^{-13}$ and the limit as $\mathsf T\to0$ to $9.5\times10^{-8}$ at $\mathsf T=10^{-9}$. On 20000 random nonnegative $2\times2$ matrices with each entry zero with probability $\frac12$, \cref{lem:evo-perron}(iii) and \cref{lem:evo-monotone} had no violation, and on 500 draws with $\mathsf b\ge1$ and $\mathsf c>0$ the matrix limit of \cref{prop:evo-transmission}(c) held to $5.1\times10^{-13}$. In \cref{ex:evo-worked} at $\mathsf T=20$ the row vector orthogonal to $r_{\mathsf M}$ is a left eigenvector of $\mathsf M$ with eigenvalue zero to rounding, so the conditional mean of that trait is lost in one generation while the heritable trait retains $0.675$; this is the exception that the statement of \cref{prop:evo-transmission}(c) records. The script reproduced every entry of \cref{tab:evo-worked}.

For the life cycle, 300 draws with $n\in\{2,\dots,6\}$ and $\mathsf T\le8$ were used with a random mean-conserving communication matrix (a symmetric weighted Laplacian plus a directed cycle, with zero column sums). The $(n+1)$-dimensional linear system for the conditional means of the members and the record was exponentiated, and the birth map was averaged over every founder set of size $\mathsf b$. The resulting map from the parent's member states and record to the daughter's mean and record equaled $\mathsf M$ applied to the parent's mean and record to $5.2\times10^{-14}$, which confirms \cref{prop:evo-transmission}(a) at the level of members without sampling. Without a record every member's coefficient equaled $\mathsf h_0/n$ to $3.5\times10^{-15}$, and with an externally written record every member's coefficient equaled $\bar{\mathsf b}e^{-\mathsf p_1\mathsf T}/n$ to $1.3\times10^{-15}$, while on 40 of the draws the forced solution computed by a fourth-order Runge-Kutta scheme for the record path $z_{\mathrm{ext}}(t)=0.3+\sin2t+0.2t$ agreed with the formula in the proof of \cref{thm:evo-record-channel}(e) to $1.4\times10^{-11}$. For all 25 pairs $(n,\mathsf b)$ with $2\le n\le6$ and two-point noises, enumeration of founder sets and noise values gave the conditional mean $\mathsf F\mathsf x$ to $2.1\times10^{-15}$, the covariance \eqref{eq:evo-noise} to $1.7\times10^{-16}$, and the variance of \cref{prop:evo-variance}(c) to $4.4\times10^{-16}$. The difference between $0.05342$ and $0.05422$ reported in \cref{chk:evo-fission} is therefore sampling error and not an error of the formula.

For selection, exact expectations over 200 random populations of 2 to 29 parents, with random nonnegative $2\times2$ transmission matrices and Poisson numbers of daughters, gave \eqref{eq:evo-price-general} for a random trait vector to $4.4\times10^{-16}$ and \eqref{eq:evo-price} to $3.3\times10^{-16}$. On 3000 draws of $(\mathsf h,\mathsf s^2,\mathsf v_{\mathrm s})$ the iterated recursion \eqref{eq:evo-balance-recursion} reached \eqref{eq:evo-balance-fixed} to $4.4\times10^{-16}$ in the variance and $2.7\times10^{-15}$ in the mean, the selection step agreed with numerical quadrature to $6.7\times10^{-16}$ and \eqref{eq:evo-balance-fitness} to $5.8\times10^{-12}$, and the equivalence of $\hat y/\mathsf y_\ast\ge\frac12$ with $\mathsf h\mathsf i\ge1-\mathsf h$ held on every draw. For 200 pairs $(\mathsf s^2,\mathsf v_{\mathrm s})$ the retained fraction was strictly increasing in $\mathsf h$ on a grid of 2001 points of $[0,1]$, and for 200 values of $\mathsf h$ in $[0.02,0.98]$ it was strictly increasing in $\mathsf s^2$ on a logarithmic grid of 1501 points of $[10^{-6},10^3]$, with value at most $3.4\times10^{-9}$ at $\mathsf s^2=10^{-12}$. The fixed point \eqref{eq:evo-directional} was reached to $2.7\times10^{-14}$, the limits of \cref{prop:evo-two-type} to $1.3\times10^{-15}$ on 971 draws with $\lvert\mathsf m\mathsf q-1\rvert>0.05$, and at $\mathsf m\mathsf q=1$ with $\mathsf m=3$ and $p_0=\frac12$ the frequency after 1000 generations was $1/2002$.

For \cref{prop:evo-bottleneck}, on 2000 draws with $n=8$ the ordering of $\mathsf h$ in $\mathsf b$ at $(\varpi_+-\varpi_-)\mathsf T=45$ agreed with the sign of \eqref{eq:evo-bottleneck-slope} on all 1981 draws with slope larger than $10^{-3}$ in absolute value. On a grid of 60 generation times log-spaced in $[0.02,60]$ no draw had a heritability that was neither nondecreasing nor nonincreasing in $\mathsf b$; this is a numerical observation on these draws and not a proved statement. Lowering the copy fidelity of the inversion instance of \cref{chk:evo-bottleneck} from $\mathsf c=1$ to $\mathsf c=0.9$ gives $\varpi_-=0.0376<\delta+\omega(1-\mathsf c\mathsf f)=0.05$, and the heritability at $\mathsf b=0$, $1$, $5$ and $10$ is then $0.89249$, $0.89312$, $0.89608$ and $0.90113$ at $\mathsf T=0.5$ and $0.43070$, $0.43092$, $0.43182$ and $0.43294$ at $\mathsf T=20$, increasing in $\mathsf b$ as the proposition predicts for long generations. With $\mathsf c=1$ the heritability of that instance was strictly decreasing in $\mathsf b$ at each of 400 generation times log-spaced in $[10^{-3},200]$.

Not tested: the covariance $\mathsf N_{\mathrm{rel}}$, which remains an unspecified input; any property of the stochastic paths beyond conditional means and the birth covariance; the sampling variance of a finite population in \cref{prop:evo-variance}(b) and \cref{prop:evo-price}; the closure \cref{ass:evo-gaussian} as a description of the fission model; and the threshold $\mathsf T_0$ of \cref{prop:evo-bottleneck}, for which no bound is given.
\end{numcheck}
